\section{Background}
\label{sec:background}

\textbf{Continual learning and parameter importance.}
Continual learning aims to reduce catastrophic forgetting as models are updated on sequential tasks~\citep{mccloskey1989catastrophic,french1999catastrophic}. Importance-based methods such as EWC~\citep{kirkpatrick2017overcoming} and SI~\citep{zenke2017continual} protect parameters judged important to earlier tasks, while gradient-based selective-update methods use calibration gradients to decide which coordinates to update~\citep{qiao2024pegp,wang2025gorp}. The latter require an additional forward--backward pass over calibration data; in our setup this adds compute and a calibration-data dependency but does not exceed the peak memory of ordinary training (Appendix~\ref{app:attribution_cost}).

\textbf{Random subspaces and intrinsic dimensionality.}
A complementary view is that the update budget may matter more than the exact coordinates selected. Random subnetworks can remain effective~\citep{frankle2019lottery,ramanujan2020whats}, fine-tuning often occupies a low-dimensional subspace~\citep{aghajanyan2020intrinsic}, and large fractions of delta parameters can be removed~\citep{yu2024dare}. Because LoRA already constrains updates to a low-rank subspace, masking imposes a second restriction within it; our question is whether gradient attribution improves over a random restriction of the same size.
